\paragraph{IV 20.11: Grassmannians and linear subvarieties (cote 25, sites 62--65).}
Section 11 begins partway down site 62 under ``Grassmaniennes'', with the old following heading ``Ratrapages'' canceled. The opening explains the need for linear subvarieties beyond the dimensions previously used and for a review of the earlier sections in this enlarged setting. It defines the functor of locally free rank-$n$ quotients of $\underline E$ and constructs its morphism to the rank-one Grassmannian, identified in the display with the projective bundle of an exterior power. The proof refers explicitly to the S\'eminaire Cartan and continues on site 63 with a closed-subscheme realization and the stated separation, quasi-compactness, projectivity and finite-presentation properties. The source's printed ``immersion ferm\'ees'', its internally spaced words, and the unresolved base letter at the start of the definition are retained. For a locally free Module, site 63 gives the affine charts indexed by choices of $n$ basis vectors, with relative dimension $n(r-n)$, and describes their corresponding subfunctors through a chosen exact sequence. Its displayed vector-bundle expression replaces a canceled tensor expression by a Hom expression. The consequences written there include local lifting of a rational point to a section and the existence of rational points in nonempty opens over an infinite field.

Sites 64--65 interpret the Grassmannian in terms of linear subvarieties $P(\underline G)$ of $P(\underline E)$, with a handwritten qualification that the general terminology imposes no fixed rank. The intrinsic characterization by smoothness, relative dimension and projective degree one is explicitly deferred to a later chapter. The polarity discussion writes a canonical comparison of the rank-$n$ and rank-$(r-n)$ Grassmannians and interprets the corresponding dual linear subvariety as the hyperplanes containing the original one. The visible unaccented $E$ in ``soit E son dual'' and the neighboring handwritten V-shaped dual accents remain as written. Site 65 records the containment criterion through the comparison of two locally free quotients, with the directions of the displayed morphisms preserved. The author then calls for a minimum collection of sorites whose full list will depend on the completed preceding material. A canceled section-12 heading is followed by the handwritten addition introducing the quotient functor with no prescribed rank and representing it by the coproduct over $n\geq0$. Its small index is partly uncertain. The concluding note expressly allows the rank, and hence relative dimension, to vary over a disconnected base. Site 66 begins the separate section on generalizing earlier results to linear sections, under the revised title ``Généralisation de résultats antérieurs aux sections linéaires''.
